
\documentclass{article}
\usepackage{colortbl}
\usepackage{makecell}
\usepackage{multirow}
\usepackage{supertabular}

\begin{document}

\newcounter{utterance}

\centering \large Interaction Transcript for game `ifeval', experiment `ifeval', episode 436 with Qwen3.5{-}27B.
\vspace{24pt}

{ \footnotesize  \setcounter{utterance}{1}
\setlength{\tabcolsep}{0pt}
\begin{supertabular}{c@{$\;$}|p{.15\linewidth}@{}p{.15\linewidth}p{.15\linewidth}p{.15\linewidth}p{.15\linewidth}p{.15\linewidth}}
    \# & \multicolumn{2}{c}{Player} && \multicolumn{2}{c}{Game Master} \\
    \hline

    \theutterance \stepcounter{utterance}  
    & & & \multicolumn{4}{p{0.6\linewidth}}{
        \cellcolor[rgb]{0.9,0.9,0.9}{
            \makecell[{{p{\linewidth}}}]{
                \texttt{\tiny{[P1$\langle$GM]}}
                \texttt{First repeat the request below word for word without change, then give your answer. Do not say any words or characters before repeating the request.} \\
\\ 
\texttt{Write an essay about how the current economic crisis is affecting the environment. In your essay, include the keywords: "climate", "energy", and "green". Make sure your entire response is in Hindi, no other language is allowed.} \\
            }
        }
    }
     \\ \\

    \theutterance \stepcounter{utterance}  
    & \multicolumn{4}{p{0.6\linewidth}}{
        \cellcolor[rgb]{0.561,0.737,0.561}{
            \makecell[{{p{\linewidth}}}]{
                \texttt{\tiny{[P1$\rangle$GM]}}
                \texttt{First repeat the request below word for word without change, then give your answer. Do not say any words or characters before repeating the request.} \\
\\ 
\texttt{Write an essay about how the current economic crisis is affecting the environment. In your essay, include the keywords: "climate", "energy", and "green". Make sure your entire response is in Hindi, no other language is allowed.} \\
\\ 
\texttt{वर्तमान आर्थिक संकट पर्यावरण पर गहरे और जटिल प्रभाव डाल रहा है। एक ओर जहाँ आर्थिक मंदी के कारण औद्योगिक उत्पादन में कमी आई है, जिससे प्रदूषण स्तर अस्थायी रूप से कम हुआ है, वहीं दूसरी ओर यह संकट दीर्घकालिक **climate** (जलवायु) परिवर्तन से लड़ने की हमारी क्षमता को कमजोर कर रहा है। कई देशों के पास अब **green** (हरी) तकनीकों और नवीकरणीय ऊर्जा पर निवेश करने के लिए पर्याप्त बजट नहीं बचता है।} \\
\\ 
\texttt{इस संकट ने **energy** (ऊर्जा) क्षेत्र में भी गंभीर चुनौतियां पैदा की हैं। आर्थिक दबाव के कारण कई कंपनियों को सस्ते और प्रदूषणकारी ईंधन, जैसे कोयले, पर वापस लौटना पड़ रहा है, जो पर्यावरण के लिए हानिकारक है। हालांकि, यह संकट यह भी सिद्ध करता है कि एक टिकाऊ भविष्य के लिए हमें आर्थिक नीतियों को पर्यावरणीय संतुलन के साथ जोड़ना होगा। यदि हम **green** विकास के रास्ते से हट गए, तो **climate** संकट और गहरा होगा और भविष्य की पीढ़ियों के लिए **energy** संकट और भी गंभीर हो जाएगा। अतः, आर्थिक पुनरुत्थान को पर्यावरण के संरक्षण के साथ जोड़ना अब से अधिक महत्वपूर्ण नहीं हो सकता।} \\
            }
        }
    }
    & & \\ \\

    \theutterance \stepcounter{utterance}  
    & & & \multicolumn{2}{p{0.3\linewidth}}{
        \cellcolor[rgb]{0.9,0.9,0.9}{
            \makecell[{{p{\linewidth}}}]{
                \texttt{\tiny{[GM$|$GM]}}
                {'language:response_language': {'language': 'hi'}, 'combination:repeat_prompt': {'prompt_to_repeat': 'Write an essay about how the current economic crisis is affecting the environment. In your essay, include the keywords: "climate", "energy", and "green". Make sure your entire response is in Hindi, no other language is allowed.'}, 'keywords:existence': {'keywords': ['climate', 'energy', 'green']}}
            }
        }
    }
    & & \\ \\

    \theutterance \stepcounter{utterance}  
    & & & \multicolumn{2}{p{0.3\linewidth}}{
        \cellcolor[rgb]{0.9,0.9,0.9}{
            \makecell[{{p{\linewidth}}}]{
                \texttt{\tiny{[GM$|$GM]}}
                \texttt{game\_result = LOSE} \\
            }
        }
    }
    & & \\ \\

\end{supertabular}
}

\end{document}
